% ===========================================================================
\section{The variance and entropy methods}\label{sec:entropy}
\warmth{0}\quad\whisper{Tensorization: variance and entropy are subadditive over product measures, so high-dim = one-dim subadditivity + a single-coordinate estimate.}

\begin{strand}
In one line: both variance (Efron–Stein) and entropy (log-Sobolev / Herbst) are \emph{subadditive}
over product measures. \keyword{Tensorize} a high-dimensional concentration into $n$ single-coordinate
contributions, each needing only a one-dimensional estimate. This route gives the tightest,
\emph{variance-aware, dimension-free} bounds, and is the analytic engine behind the geometry of
§\ref{sec:geometry} and Talagrand (§\ref{sec:talagrand}).
\end{strand}

\block{The paradigm}
\trigger{You want the tightest \emph{variance-aware} bound, or to prove a function sub-Gaussian without a ``sum / martingale'' structure.}
\keyword{Herbst's argument}: if the entropy obeys $\Ent[e^{\lambda f}]\le \tfrac{\lambda^2\sigma^2}{2}\,\E[e^{\lambda f}]$
for all $\lambda$, then $f-\E f$ is $\sigma$-sub-Gaussian. Combine with the \keyword{tensorization} of entropy
$\Ent[g]\le \sum_i \E\,\Ent_i[g]$ to reduce the problem to a single coordinate.

\begin{strand}
Three-line heuristic: (1) \keyword{Efron–Stein} is the tensorization of variance, giving the \emph{free}
(no i.i.d.\ needed) bound $\Var f\le \tfrac12\sum_i\E(f-f^{(i)})^2$; (2) replace $\Var$ by $\Ent$ and ``square''
by ``exponential'': tensorization + a one-dimensional log-Sobolev $=$ exponential concentration (Herbst);
(3) the Gaussian measure satisfies a \emph{dimension-free} log-Sobolev inequality (Gross), the algebraic source
of the Gaussian concentration in §\ref{sec:geometry}.
\end{strand}

\block{Cheat-sheet: variance and entropy}
\noindent\begin{tabularx}{\linewidth}{@{}l l L@{}}
\toprule
{\color{indigo}Theorem} & {\color{indigo}Setting} & \multicolumn{1}{l}{\color{indigo}Gives / when to use}\\
\midrule
Efron–Stein     & product measure, $f^{(i)}=$ resample coord.\ $i$ & $\Var f\le\tfrac12\sum_i\E(f-f^{(i)})^2$; the ``bounded difference'' for variance\\
Bounded diff.\ via ES & bounded diff.\ $c_i$               & $\Var f\le\tfrac14\sum c_i^2$; variance from the martingale family\\
Herbst argument & $\Ent[e^{\lambda f}]\le\tfrac{\lambda^2\sigma^2}2\E e^{\lambda f}$ & $\Rightarrow f$ is $\sigma$-sub-Gaussian\\
Gaussian LSI    & $\gamma=N(0,I_n)$               & $\Ent[g^2]\le 2\E\|\nabla g\|^2$; dimension-free\\
Modified LSI    & bounded / Bernoulli coords      & tensorize + single-coord LSI $\Rightarrow$ concentration\\
\bottomrule
\end{tabularx}

\block{Main results \& proof skeletons}

\begin{theorem}[Efron–Stein inequality]\label{thm:efron-stein}
$X_1,\dots,X_n$ independent, $X_i'$ an i.i.d.\ copy, $f^{(i)}$ replacing coordinate $i$ by $X_i'$. Then
$\displaystyle \Var f\le \tfrac12\sum_{i=1}^n\E\big[(f-f^{(i)})^2\big]$.
\end{theorem}
\begin{proof}
Skeleton: use the orthogonal martingale-difference decomposition $\Var f=\sum_i\E[\Delta_i^2]$ with
$\Delta_i=\E[f\mid X_{1:i}]-\E[f\mid X_{1:i-1}]$; then Jensen to replace $\Delta_i$ by the symmetrized $(f-f^{(i)})$.
\TODO{fill the orthogonal decomposition + symmetrization, two steps}
\end{proof}

\begin{theorem}[Herbst's argument]\label{thm:herbst}
If $\Ent[e^{\lambda f}]\le\tfrac{\lambda^2\sigma^2}2\E[e^{\lambda f}]$ for all $\lambda$, then
$\E[e^{\lambda(f-\E f)}]\le e^{\lambda^2\sigma^2/2}$, so $f$ is $\sigma$-sub-Gaussian.
\end{theorem}
\begin{proof}
Skeleton: let $H(\lambda)=\tfrac1\lambda\log\E e^{\lambda f}$; the entropy inequality is exactly the rewriting
$H'(\lambda)\le\sigma^2/2$; integrate to get $H(\lambda)-H(0^+)\le\lambda\sigma^2/2$, with $H(0^+)=\E f$.
\TODO{fill the differential identity $\Ent[e^{\lambda f}]=\lambda^2\E e^{\lambda f}\,H'(\lambda)$}
\end{proof}

\begin{theorem}[Gaussian log-Sobolev, Gross]\label{thm:gaussian-lsi}
For $\gamma=N(0,I_n)$ and smooth $g$, $\Ent_\gamma[g^2]\le 2\,\E_\gamma\|\nabla g\|^2$.
\end{theorem}
\begin{proof}\TODO{fill: the one-dimensional LSI (Ornstein–Uhlenbeck semigroup, or CLT limit from two-point Bernoulli) + tensorization to $n$ dimensions}\end{proof}

\begin{strand}
Why it goes this way: Gaussian LSI $+$ Herbst $\Rightarrow$ ``Gaussian concentration $e^{-t^2/2L^2}$ for an
$L$-Lipschitz function, dimension-free''. This is the algebraic path that \emph{computes} the geometric
statement of §\ref{sec:geometry}: curvature (the LSI constant) turns directly into the tail rate.
\end{strand}

\loose{Four faces: concentration $\Leftrightarrow$ isoperimetry $\Leftrightarrow$ log-Sobolev $\Leftrightarrow$ transportation (Marton). Herbst is the ``LSI $\Rightarrow$ concentration'' face; the other three are to be filled, see §\ref{sec:geometry} and §9 open threads.}%
\pick{tensorize}\quad Tensorization reappears in §\ref{sec:talagrand} in Marton's transportation form.
